% id: 149-10
% book: 베이직쎈 미적분1 · 09 정적분의 활용
% section: 학교 시험 기출로 실전 감각 UP!
% figure: tikz:fig-149-10
% answer: 
% answer_source: 
% variant_level: 0 (원본 전사)
% 이 파일은 build.mjs 가 items.json 에서 만든다. 고칠 때는 items.json 을 고친다.
\smash{\raisebox{1.5mm}{\dmkichul{베이직쎈 미적분1 149쪽 10번 · 꼭 나와요}}}
\noindent\begin{minipage}[t]{0.58\linewidth}\vspace{0pt}\raggedright 원점을 출발하여 수직선 위를 움직이는 점 $\pt{P}$의 시각 $t$에서의 속도 $v(t)$의 그래프가 오른쪽 그림과 같다. 점 $\pt{P}$가 출발 후 다시 원점을 지날 때의 시각은?\end{minipage}\hfill\begin{minipage}[t]{0.4\linewidth}\vspace{0pt}\centering \resizebox{\ifdim\width>\linewidth\linewidth\else\width\fi}{!}{% 149-10(149쪽) — 원점을 출발한 점 P 의 시각 t 에서의 속도 v(t) 의 그래프(꺾은선 (0,0)-(2,2)-(4,-2)-(5,-2)-(6,0)). 스캔 화질이 낮아 TikZ 로 다시 그림.
% 원문에 있는 것만 옮긴다: 꺾은선 · 안내 점선 다섯 줄 · 라벨 v(t), t, O, 2·$-2$(세로축), 2(가로축 아래), 3 4 5 6(가로축 위). 원문에 없는 값은 적지 않는다.
\begin{tikzpicture}[x=0.36cm, y=0.36cm, line width=0.6pt]
  \draw[->, >=stealth, line width=0.7pt] (-0.8,0) -- (6.9,0) node[below=1pt, font=\small] {$t$};
  \draw[->, >=stealth, line width=0.7pt] (0,-3.3) -- (0,3.1) node[left=1pt, font=\small] {$v(t)$};
  \draw[densely dashed, line width=0.4pt] (0,2) -- (2,2);
  \draw[densely dashed, line width=0.4pt] (2,2) -- (2,0);
  \draw[densely dashed, line width=0.4pt] (0,-2) -- (4,-2);
  \draw[densely dashed, line width=0.4pt] (4,0) -- (4,-2);
  \draw[densely dashed, line width=0.4pt] (5,0) -- (5,-2);
  \draw (0,0) -- (2,2) -- (4,-2) -- (5,-2) -- (6,0);
  \node[left=1pt, font=\small] at (0,2) {$2$};
  \node[left=1pt, font=\small] at (0,-2) {$-2$};
  \node[below left=-2pt and -2.5pt, font=\small] at (0,0) {$\pt{O}$};
  \node[below=1pt, font=\small] at (2,0) {$2$};
  \node[above=1pt, font=\small] at (3.3,0) {$3$};
  \node[above=1pt, font=\small] at (4.05,0) {$4$};
  \node[above=1pt, font=\small] at (5,0) {$5$};
  \node[above=1pt, font=\small] at (6,0) {$6$};
\end{tikzpicture}
}\end{minipage}\par
\dmchoicesauto{}{$2$}{$3$}{$4$}{$5$}{$6$}
